\scanpage{49}
\margin{\unsure{84|48}}
\begin{flushright}6\end{flushright}

Continuing in this way we construct the minimal model
\[
\Lambda(x_1,x_2,\ldots:d)\underset{\rho}{\longrightarrow}\mathcal E
\]
satisfying properties i) and ii) above.

\uline{Theorem B}: Let $X$ be a triangulated space which is simply connected. If
\[
\Lambda(x_1,x_2,\ldots)\underset{\rho}{\longrightarrow}\mathcal E
\]
is a minimal model of the de Rham algebra then

i) the rank of $\pi_i X$ is the number of variables in degree $i$.

ii) the rank of the Hurewicz homomorphism
$\pi_i X\longrightarrow H_i X$ is the number of $d$-closed generators in degree $i$.

iii) the Whitehead products on homotopy are described by the quadratic terms of the $d$-formulae of generators.

Theorem B is proved by a not so standard Postnikov argument using Theorem A and the commutativity of the wedge product multiplication in $\mathcal E$. The argument is based on the beautifully simple Guy Hirsh method of computing the cohomology of a principal $G$-bundle over a manifold using differential forms (see appendix H). The Hirsch method was suggested by Rene Thom in the study of Postnikov systems in the Cartan Seminar 1954.

The proof of theorem B shows the model can be used to directly construct the \uline{rational} Postnikov system
